\documentclass[12pt,a4paper]{article}
\usepackage[UTF8]{ctex}
\usepackage{geometry}
\usepackage{listings}
\usepackage{xcolor}
\usepackage{hyperref}
\usepackage{booktabs}
\usepackage{float}
\usepackage{tcolorbox}
\usepackage{enumitem}

\geometry{left=2.5cm,right=2.5cm,top=2.5cm,bottom=2.5cm}

\lstset{
    language=Python,
    basicstyle=\ttfamily\small,
    keywordstyle=\color{blue}\bfseries,
    commentstyle=\color{gray},
    stringstyle=\color{orange},
    numbers=left,
    numberstyle=\tiny\color{gray},
    stepnumber=1,
    numbersep=8pt,
    backgroundcolor=\color{gray!10},
    frame=single,
    frameround=fttt,
    breaklines=true,
    showstringspaces=false,
    tabsize=4,
    captionpos=b,
    xleftmargin=2em,
    xrightmargin=1em,
    aboveskip=1em,
    belowskip=1em
}

\tcbuselibrary{skins,breakable}
\newtcolorbox{knowledgebox}[1][]{
    colback=blue!5!white,
    colframe=blue!50!black,
    fonttitle=\bfseries,
    title=#1,
    breakable
}

\newtcolorbox{practicebox}[1][]{
    colback=green!5!white,
    colframe=green!50!black,
    fonttitle=\bfseries,
    title=#1,
    breakable
}

\newtcolorbox{tipbox}[1][]{
    colback=yellow!10!white,
    colframe=orange!70!black,
    fonttitle=\bfseries,
    title=#1,
    breakable
}

\hypersetup{
    colorlinks=true,
    linkcolor=blue,
    filecolor=magenta,      
    urlcolor=cyan,
    bookmarks=true,
    pdfstartview=FitH
}

\title{\textbf{Pandas数据处理实战手册}\\[0.5em]\large 基于真实数据案例}
\author{竞赛培训组}
\date{\today}

\begin{document}

\maketitle
\tableofcontents
\newpage

%========================================
\section{四个练习数据文件介绍}
%========================================

\subsection{数据文件概览}

本手册所有示例代码都围绕以下\textbf{四个CSV文件}展开，请确保这些文件和代码在同一目录下：

\begin{table}[H]
\centering
\caption{练习数据文件说明}
\begin{tabular}{|l|l|l|l|}
\hline
\textbf{文件名} & \textbf{数据量} & \textbf{主题} & \textbf{主要问题} \\
\hline
01\_电商订单数据\_练习.csv & 50行 & 电商订单 & 重复、缺失、异常值、格式混乱 \\
\hline
02\_传感器数据\_练习.csv & 144行 & 传感器读数 & 时间序列、格式问题 \\
\hline
03\_缺失值处理\_练习.csv & 30行 & 员工信息 & 各种缺失值场景 \\
\hline
04\_数据格式清洗\_练习.csv & 30行 & 用户信息 & 格式不统一、单位混乱 \\
\hline
\end{tabular}
\end{table}

\subsection{01\_电商订单数据\_练习.csv}

\begin{knowledgebox}[数据字段说明]
\begin{itemize}
    \item \textbf{order\_id}：订单编号
    \item \textbf{user\_id}：用户ID（部分缺失）
    \item \textbf{product\_name}：商品名称
    \item \textbf{price}：价格（带"元"单位）
    \item \textbf{quantity}：数量（有负数异常值）
    \item \textbf{order\_date}：订单日期（多种格式）
    \item \textbf{address}：地址
    \item \textbf{phone}：手机号（带横线）
    \item \textbf{tags}：商品标签（JSON格式）
\end{itemize}
\end{knowledgebox}

\subsection{02\_传感器数据\_练习.csv}

\begin{knowledgebox}[数据字段说明]
\begin{itemize}
    \item \textbf{timestamp}：时间戳
    \item \textbf{sensor\_id}：传感器ID
    \item \textbf{temperature}：温度
    \item \textbf{humidity}：湿度（部分缺失）
    \item \textbf{pressure}：压力
    \item \textbf{status}：设备状态
\end{itemize}
\end{knowledgebox}

\subsection{03\_缺失值处理\_练习.csv}

\begin{knowledgebox}[数据字段说明]
\begin{itemize}
    \item \textbf{id}：编号
    \item \textbf{name}：姓名
    \item \textbf{age}：年龄（部分缺失）
    \item \textbf{gender}：性别（部分缺失）
    \item \textbf{salary}：工资（部分缺失）
    \item \textbf{department}：部门（部分缺失）
    \item \textbf{join\_date}：入职日期（部分缺失）
    \item \textbf{phone}：电话
    \item \textbf{email}：邮箱
    \item \textbf{score}：评分（部分缺失）
\end{itemize}
\end{knowledgebox}

\subsection{04\_数据格式清洗\_练习.csv}

\begin{knowledgebox}[数据字段说明]
\begin{itemize}
    \item \textbf{id}：编号
    \item \textbf{name}：姓名（带空格）
    \item \textbf{birth\_date}：出生日期（多种格式）
    \item \textbf{height}：身高（带单位）
    \item \textbf{weight}：体重（带单位）
    \item \textbf{income}：收入（带k单位）
    \item \textbf{phone}：电话
    \item \textbf{address}：地址
    \item \textbf{product\_tags}：商品标签
\end{itemize}
\end{knowledgebox}

%========================================
\section{Python与Pandas基础}
%========================================

\subsection{安装与导入}

\begin{lstlisting}[caption={安装Pandas}]
# 在命令行中运行
pip install pandas
\end{lstlisting}

\begin{lstlisting}[caption={导入Pandas}]
import pandas as pd
import numpy as np
\end{lstlisting}

\subsection{读取CSV文件}

\begin{knowledgebox}[知识点：pd.read\_csv()]
这是Pandas中最常用的函数，用于读取CSV文件。
\end{knowledgebox}

\begin{lstlisting}[caption={读取四个练习文件}]
# 读取电商订单数据
df\_order = pd.read\_csv('01\_电商订单数据\_练习.csv')

# 读取传感器数据
df\_sensor = pd.read\_csv('02\_传感器数据\_练习.csv')

# 读取缺失值练习数据
df\_missing = pd.read\_csv('03\_缺失值处理\_练习.csv')

# 读取格式清洗练习数据
df\_messy = pd.read\_csv('04\_数据格式清洗\_练习.csv')
\end{lstlisting}

\begin{practicebox}[动手练习1：读取并查看数据]
\textbf{任务：}读取四个CSV文件，查看每个文件的基本信息。

\begin{lstlisting}
import pandas as pd

# 读取四个文件
df\_order = pd.read\_csv('01\_电商订单数据\_练习.csv')
df\_sensor = pd.read\_csv('02\_传感器数据\_练习.csv')
df\_missing = pd.read\_csv('03\_缺失值处理\_练习.csv')
df\_messy = pd.read\_csv('04\_数据格式清洗\_练习.csv')

# 查看电商订单数据的形状
print("电商订单数据形状：", df\_order.shape)
print("\n前5行：")
print(df\_order.head())
\end{lstlisting}
\end{practicebox}

%========================================
\section{数据查看与探索}
%========================================

\subsection{查看数据基本信息}

\begin{lstlisting}[caption={查看数据形状和列名}]
# 查看数据形状（行数，列数）
print(df\_order.shape)

# 查看列名
print(df\_order.columns)

# 查看数据类型
print(df\_order.dtypes)
\end{lstlisting}

\subsection{查看数据内容}

\begin{lstlisting}[caption={查看数据内容}]
# 查看前5行
print(df\_order.head())

# 查看后5行
print(df\_order.tail())

# 查看前10行
print(df\_order.head(10))
\end{lstlisting}

\subsection{统计信息}

\begin{lstlisting}[caption={使用describe()查看统计信息}]
# 查看数值列的统计信息
print(df\_order.describe())

# 查看缺失值统计
print(df\_order.isnull().sum())
\end{lstlisting}

\begin{practicebox}[动手练习2：探索传感器数据]
\textbf{任务：}使用02\_传感器数据\_练习.csv，查看数据的基本统计信息。

\begin{lstlisting}
import pandas as pd

df\_sensor = pd.read\_csv('02\_传感器数据\_练习.csv')

# 查看形状
print("数据形状：", df\_sensor.shape)

# 查看统计信息
print("\n统计摘要：")
print(df\_sensor.describe())

# 查看有多少个不同的传感器
print("\n传感器数量：", df\_sensor['sensor\_id'].nunique())
print("传感器列表：", df\_sensor['sensor\_id'].unique())
\end{lstlisting}
\end{practicebox}

%========================================
\section{数据选择}
%========================================

\subsection{选择列}

\begin{lstlisting}[caption={选择单列和多列}]
# 选择单列（商品名称）
products = df\_order['product\_name']

# 选择多列（订单ID、商品名称、价格）
order\_info = df\_order[['order\_id', 'product\_name', 'price']]
\end{lstlisting}

\subsection{选择行}

\begin{lstlisting}[caption={使用iloc选择行}]
# 选择第1行（索引为0）
row = df\_order.iloc[0]

# 选择前5行
rows = df\_order.iloc[0:5]

# 选择第10到20行
rows = df\_order.iloc[10:20]
\end{lstlisting}

\subsection{条件筛选}

\begin{lstlisting}[caption={条件筛选示例}]
# 筛选数量大于1的订单
large\_orders = df\_order[df\_order['quantity'] > 1]

# 筛选特定商品
iphone\_orders = df\_order[df\_order['product\_name'] == 'iPhone 14 Pro']
\end{lstlisting}

\begin{practicebox}[动手练习3：筛选缺失值数据]
\textbf{任务：}使用03\_缺失值处理\_练习.csv，筛选出年龄缺失的员工。

\begin{lstlisting}
import pandas as pd

df\_missing = pd.read\_csv('03\_缺失值处理\_练习.csv')

# 筛选年龄缺失的行
missing\_age = df\_missing[df\_missing['age'].isnull()]

print("年龄缺失的员工：")
print(missing\_age[['name', 'age', 'department']])

print(f"\n共 {len(missing\_age)} 人年龄缺失")
\end{lstlisting}
\end{practicebox}

%========================================
\section{数据清洗：处理缺失值}
%========================================

\subsection{检测缺失值}

\begin{lstlisting}[caption={检测缺失值}]
# 统计每列的缺失值数量
print(df\_missing.isnull().sum())

# 计算缺失值比例
missing\_ratio = df\_missing.isnull().sum() / len(df\_missing) * 100
print(missing\_ratio)
\end{lstlisting}

\subsection{填充缺失值}

\begin{lstlisting}[caption={填充缺失值的方法}]
# 用固定值填充
df\_missing['gender'] = df\_missing['gender'].fillna('未知')

# 用平均值填充（工资列）
mean\_salary = df\_missing['salary'].mean()
df\_missing['salary'] = df\_missing['salary'].fillna(mean\_salary)

# 用中位数填充（年龄列）
median\_age = df\_missing['age'].median()
df\_missing['age'] = df\_missing['age'].fillna(median\_age)

# 用众数填充（部门列）
mode\_dept = df\_missing['department'].mode()[0]
df\_missing['department'] = df\_missing['department'].fillna(mode\_dept)
\end{lstlisting}

\begin{practicebox}[动手练习4：处理缺失值]
\textbf{任务：}使用03\_缺失值处理\_练习.csv，完成缺失值处理。

\begin{lstlisting}
import pandas as pd

df = pd.read\_csv('03\_缺失值处理\_练习.csv')

print("处理前缺失值统计：")
print(df.isnull().sum())

# 用中位数填充年龄
median\_age = df['age'].median()
df['age'] = df['age'].fillna(median\_age)

# 用众数填充性别
mode\_gender = df['gender'].mode()[0]
df['gender'] = df['gender'].fillna(mode\_gender)

# 用平均值填充工资
mean\_salary = df['salary'].mean()
df['salary'] = df['salary'].fillna(mean\_salary)

print("\n处理后缺失值统计：")
print(df.isnull().sum())
\end{lstlisting}
\end{practicebox}

%========================================
\section{数据清洗：处理重复值}
%========================================

\subsection{检测和删除重复值}

\begin{lstlisting}[caption={处理重复值}]
# 检测重复值
duplicates = df\_order.duplicated()
print(f"重复行数：{duplicates.sum()}")

# 查看重复的行
duplicate\_rows = df\_order[df\_order.duplicated(keep=False)]
print(duplicate\_rows)

# 删除重复值（保留第一个）
df\_order\_clean = df\_order.drop\_duplicates()
\end{lstlisting}

\begin{practicebox}[动手练习5：处理电商订单重复值]
\textbf{任务：}使用01\_电商订单数据\_练习.csv，检测并删除重复订单。

\begin{lstlisting}
import pandas as pd

df = pd.read\_csv('01\_电商订单数据\_练习.csv')

print(f"原始数据行数：{len(df)}")

# 检测重复值
n\_duplicates = df.duplicated().sum()
print(f"重复行数：{n\_duplicates}")

# 查看重复的行
duplicate\_rows = df[df.duplicated(keep=False)]
print("\n重复的行：")
print(duplicate\_rows[['order\_id', 'product\_name', 'price']])

# 删除重复值
df\_clean = df.drop\_duplicates()
print(f"\n删除后行数：{len(df\_clean)}")
\end{lstlisting}
\end{practicebox}

%========================================
\section{数据清洗：处理异常值}
%========================================

\subsection{检测异常值}

\begin{lstlisting}[caption={使用IQR方法检测异常值}]
# 计算IQR
Q1 = df\_order['quantity'].quantile(0.25)
Q3 = df\_order['quantity'].quantile(0.75)
IQR = Q3 - Q1

# 定义异常值范围
lower\_bound = Q1 - 1.5 * IQR
upper\_bound = Q3 + 1.5 * IQR

# 找出异常值
outliers = df\_order[(df\_order['quantity'] < lower\_bound) | 
                    (df\_order['quantity'] > upper\_bound)]
print(f"异常值数量：{len(outliers)}")
\end{lstlisting}

\subsection{处理异常值}

\begin{lstlisting}[caption={处理异常值}]
# 将负数数量修正为1
df\_order.loc[df\_order['quantity'] < 0, 'quantity'] = 1
\end{lstlisting}

%========================================
\section{数据格式转换}
%========================================

\subsection{字符串处理}

\begin{lstlisting}[caption={处理价格列}]
# 查看价格列的格式
print(df\_order['price'].head())

# 去除"元"字，转换为数值
df\_order['price\_num'] = df\_order['price'].str.replace('元', '')
df\_order['price\_num'] = df\_order['price\_num'].astype(float)
\end{lstlisting}

\subsection{处理手机号}

\begin{lstlisting}[caption={标准化手机号}]
# 去除横线
df\_order['phone\_clean'] = df\_order['phone'].str.replace('-', '')
\end{lstlisting}

\begin{practicebox}[动手练习6：清洗格式混乱数据]
\textbf{任务：}使用04\_数据格式清洗\_练习.csv，清洗姓名和收入列。

\begin{lstlisting}
import pandas as pd

df = pd.read\_csv('04\_数据格式清洗\_练习.csv')

# 查看原始数据
print("姓名处理前：")
print(df['name'].head())

# 去除姓名中的空格
df['name\_clean'] = df['name'].str.strip()

print("\n姓名处理后：")
print(df['name\_clean'].head())

# 处理收入（将12k转换为12000）
def convert\_income(income\_str):
    income\_str = str(income\_str).strip()
    if 'k' in income\_str.lower():
        number = float(income\_str.lower().replace('k', ''))
        return number * 1000
    else:
        return float(income\_str)

df['income\_num'] = df['income'].apply(convert\_income)

print("\n收入转换：")
print(df[['income', 'income\_num']].head(10))
\end{lstlisting}
\end{practicebox}

%========================================
\section{数据分组与统计}
%========================================

\subsection{分组统计}

\begin{lstlisting}[caption={按部门统计工资}]
# 按部门分组，计算平均工资
dept\_avg = df\_missing.groupby('department')['salary'].mean()
print(dept\_avg)

# 多列分组统计
stats = df\_missing.groupby('department').agg({
    'salary': ['mean', 'max', 'min'],
    'age': 'mean'
})
print(stats)
\end{lstlisting}

\begin{practicebox}[动手练习7：统计电商数据]
\textbf{任务：}使用01\_电商订单数据\_练习.csv，统计各商品的总销量。

\begin{lstlisting}
import pandas as pd

df = pd.read\_csv('01\_电商订单数据\_练习.csv')

# 处理价格（去除"元"字）
df['price\_num'] = df['price'].str.replace('元', '').astype(float)

# 计算订单金额
df['total\_amount'] = df['price\_num'] * df['quantity']

# 按商品分组统计
product\_stats = df.groupby('product\_name').agg({
    'quantity': 'sum',
    'total\_amount': 'sum'
}).sort\_values('total\_amount', ascending=False)

print("商品销售统计TOP10：")
print(product\_stats.head(10))
\end{lstlisting}
\end{practicebox}

%========================================
\section{综合实战：完整数据处理流程}
%========================================

\subsection{案例：完整的电商订单数据处理}

\begin{lstlisting}[caption={完整的数据处理流程}]
import pandas as pd
import numpy as np

# 第1步：读取数据
df = pd.read\_csv('01\_电商订单数据\_练习.csv')
print(f"原始数据：{len(df)}行")

# 第2步：删除重复值
df = df.drop\_duplicates()
print(f"删除重复后：{len(df)}行")

# 第3步：处理缺失值
df['user\_id'] = df['user\_id'].fillna('UNKNOWN')

# 第4步：处理异常值（负数数量）
df.loc[df['quantity'] < 0, 'quantity'] = 1

# 第5步：格式转换（价格）
df['price\_num'] = df['price'].str.replace('元', '').astype(float)

# 第6步：格式转换（手机号）
df['phone\_clean'] = df['phone'].str.replace('-', '')

# 第7步：计算订单金额
df['total\_amount'] = df['price\_num'] * df['quantity']

# 第8步：统计分析
product\_sales = df.groupby('product\_name')['total\_amount'].sum()
print("\n各商品销售额：")
print(product\_sales.sort\_values(ascending=False).head())

# 第9步：保存结果
df.to\_csv('清洗后的订单数据.csv', index=False)
print("\n数据已保存到：清洗后的订单数据.csv")
\end{lstlisting}

%========================================
\section{竞赛模拟题}
%========================================

\subsection{题目一：电商订单数据清洗}

\begin{practicebox}[竞赛模拟题1]
\textbf{题目要求：}

使用\textbf{01\_电商订单数据\_练习.csv}，完成以下任务：

\begin{enumerate}
    \item 读取数据，查看基本信息（形状、列名、前5行）
    \item 统计缺失值数量
    \item 删除重复订单
    \item 处理user\_id缺失值（用'UNKNOWN'填充）
    \item 处理负数数量（修正为1）
    \item 提取价格中的数字，转换为数值型
    \item 标准化手机号格式（去除横线）
    \item 计算每个订单的总金额（价格×数量）
    \item 统计每个商品的总销售额和总销量
    \item 保存清洗后的数据到新的CSV文件
\end{enumerate}

\textbf{参考答案：}
\begin{lstlisting}
import pandas as pd

# 读取数据
df = pd.read\_csv('01\_电商订单数据\_练习.csv')

# 查看基本信息
print(f"数据形状：{df.shape}")
print(f"列名：{df.columns.tolist()}")
print(df.head())

# 统计缺失值
print("\n缺失值统计：")
print(df.isnull().sum())

# 删除重复值
df = df.drop\_duplicates()

# 填充缺失值
df['user\_id'] = df['user\_id'].fillna('UNKNOWN')

# 处理异常值
df.loc[df['quantity'] < 0, 'quantity'] = 1

# 格式转换
df['price\_num'] = df['price'].str.replace('元', '').astype(float)
df['phone\_clean'] = df['phone'].str.replace('-', '')

# 计算金额
df['total\_amount'] = df['price\_num'] * df['quantity']

# 统计分析
product\_stats = df.groupby('product\_name').agg({
    'total\_amount': 'sum',
    'quantity': 'sum'
}).sort\_values('total\_amount', ascending=False)

print("\n商品销售统计：")
print(product\_stats.head(10))

# 保存结果
df.to\_csv('清洗后的订单数据.csv', index=False)
\end{lstlisting}
\end{practicebox}

\subsection{题目二：员工信息数据处理}

\begin{practicebox}[竞赛模拟题2]
\textbf{题目要求：}

使用\textbf{03\_缺失值处理\_练习.csv}，完成以下任务：

\begin{enumerate}
    \item 读取数据，查看缺失值情况
    \item 用中位数填充age列的缺失值
    \item 用平均值填充salary列的缺失值
    \item 用众数填充department列的缺失值
    \item 用众数填充gender列的缺失值
    \item 删除score列仍有缺失值的行
    \item 按department分组，统计平均工资和平均年龄
    \item 保存处理后的数据
\end{enumerate}

\textbf{参考答案：}
\begin{lstlisting}
import pandas as pd

# 读取数据
df = pd.read\_csv('03\_缺失值处理\_练习.csv')

# 查看缺失值
print("缺失值统计：")
print(df.isnull().sum())

# 用中位数填充年龄
median\_age = df['age'].median()
df['age'] = df['age'].fillna(median\_age)

# 用平均值填充工资
mean\_salary = df['salary'].mean()
df['salary'] = df['salary'].fillna(mean\_salary)

# 用众数填充部门和性别
mode\_dept = df['department'].mode()[0]
df['department'] = df['department'].fillna(mode\_dept)

mode\_gender = df['gender'].mode()[0]
df['gender'] = df['gender'].fillna(mode\_gender)

# 删除score缺失的行
df = df.dropna(subset=['score'])

# 按部门统计
dept\_stats = df.groupby('department').agg({
    'salary': 'mean',
    'age': 'mean'
}).round(2)

print("\n部门统计：")
print(dept\_stats)

# 保存结果
df.to\_csv('处理后的员工数据.csv', index=False)
\end{lstlisting}
\end{practicebox}

%========================================
\section{代码速查表}
%========================================

\subsection{文件操作}

\begin{lstlisting}[caption={读取和保存CSV}]
# 读取CSV
df = pd.read\_csv('文件名.csv')

# 保存CSV
df.to\_csv('新文件名.csv', index=False)
\end{lstlisting}

\subsection{数据查看}

\begin{lstlisting}[caption={数据查看}]
df.shape          # 形状（行数，列数）
df.columns        # 列名
df.head()         # 前5行
df.tail()         # 后5行
df.info()         # 详细信息
df.describe()     # 统计摘要
df.dtypes         # 数据类型
\end{lstlisting}

\subsection{数据选择}

\begin{lstlisting}[caption={数据选择}]
# 选择列
df['列名']
df[['列1', '列2']]

# 选择行
df.iloc[0]        # 第1行
df.iloc[0:5]      # 前5行

# 条件筛选
df[df['列名'] > 值]
\end{lstlisting}

\subsection{数据清洗}

\begin{lstlisting}[caption={数据清洗}]
# 缺失值
df.isnull().sum()                 # 统计缺失值
df['列'].fillna(值)              # 填充缺失值
df.dropna()                       # 删除缺失值

# 重复值
df.duplicated().sum()             # 统计重复值
df.drop\_duplicates()              # 删除重复值

# 异常值（IQR方法）
Q1 = df['列'].quantile(0.25)
Q3 = df['列'].quantile(0.75)
IQR = Q3 - Q1
outliers = df[(df['列'] < Q1-1.5*IQR) | (df['列'] > Q3+1.5*IQR)]
\end{lstlisting}

\subsection{数据转换}

\begin{lstlisting}[caption={数据转换}]
# 类型转换
df['列'].astype(float)
df['列'].astype(str)

# 字符串处理
df['列'].str.replace('旧', '新')
df['列'].str.strip()

# 应用函数
df['新列'] = df['列'].apply(函数)
\end{lstlisting}

\subsection{分组统计}

\begin{lstlisting}[caption={分组统计}]
# 简单分组
df.groupby('分组列')['统计列'].mean()

# 多统计量
df.groupby('分组列').agg({
    '列1': 'mean',
    '列2': ['mean', 'sum']
})
\end{lstlisting}

\end{document}
